\documentclass[a4paper,11pt]{ctexart}
\usepackage{amsmath, amssymb, amsfonts}
\usepackage{geometry}
\geometry{left=1.5cm, right=1.5cm, top=2.0cm, bottom=2.0cm, headsep=0.3cm, footskip=1cm}
\usepackage{listings}
\usepackage{xcolor}
\usepackage{tikz}
\usepackage{pgfplots}
\usepackage{hyperref}
\usepackage{fancyhdr}
\usepackage{tcolorbox}
\tcbuselibrary{breakable, skins}
\usepackage{array}
\usepackage{booktabs}
\usepackage[shortlabels]{enumitem}
\usepackage{needspace}
\usepackage{multirow}

\AtBeginDocument{
  \hypersetup{colorlinks=true, linkcolor=blue, filecolor=magenta, urlcolor=cyan}
}

\setlist{nosep, itemsep=0.8ex, parsep=0.1em}
\setlist[itemize]{leftmargin=1.8em}
\setlist[enumerate]{leftmargin=1.8em}

\definecolor{codegreen}{rgb}{0,0.6,0}
\definecolor{codegray}{rgb}{0.5,0.5,0.5}
\definecolor{codepurple}{rgb}{0.58,0,0.82}
\definecolor{codebg}{rgb}{0.98,0.98,0.98}
\definecolor{tipsblue}{rgb}{0.85,0.93,1.0}
\definecolor{highlightyellow}{rgb}{1.0,0.97,0.8}

\lstdefinestyle{mystyle}{
    backgroundcolor=\color{codebg},
    commentstyle=\color{codegreen},
    keywordstyle=\color{magenta}\bfseries,
    numberstyle=\tiny\color{codegray},
    stringstyle=\color{codepurple},
    basicstyle=\ttfamily\small,
    breakatwhitespace=false, breaklines=true,
    captionpos=b, keepspaces=true,
    numbers=left, numbersep=6pt,
    showspaces=false, showstringspaces=false,
    showtabs=false, tabsize=2,
    language=Python, frame=single,
    framerule=0.5pt, rulecolor=\color{gray!40}
}
\lstset{style=mystyle}

\newtcolorbox{problembox}[1][]{colback=white,colframe=blue!60!black,title=\textbf{#1},fonttitle=\bfseries\small,boxrule=1pt,arc=3pt,outer arc=3pt,coltitle=black,before skip=10pt,after skip=8pt}
\newtcolorbox{solutionbox}[1][]{enhanced,breakable,colback=green!3!white,colframe=green!50!black,title=\textbf{#1},fonttitle=\bfseries\small,boxrule=1pt,arc=3pt,outer arc=3pt,coltitle=black,before skip=8pt,after skip=10pt}
\newtcolorbox{metaphorbox}[1][]{colback=orange!5!white,colframe=orange!70!black,title=\textbf{#1},fonttitle=\bfseries\small,boxrule=1pt,arc=3pt,outer arc=3pt,coltitle=black,before skip=8pt,after skip=8pt}
\newtcolorbox{beginnerbox}[1][]{colback=tipsblue,colframe=blue!50!black,title=\textbf{#1},fonttitle=\bfseries\small,boxrule=1pt,arc=3pt,outer arc=3pt,coltitle=black,before skip=8pt,after skip=8pt}
\newtcolorbox{appbox}[1][]{colback=green!3!white,colframe=green!60!black,title=\textbf{#1},fonttitle=\bfseries\small,boxrule=1pt,arc=3pt,outer arc=3pt,coltitle=black,before skip=8pt,after skip=10pt}
\newtcolorbox{keybox}[1][]{colback=highlightyellow,colframe=orange!80!black,title=\textbf{#1},fonttitle=\bfseries\small,boxrule=1pt,arc=3pt,outer arc=3pt,coltitle=black,before skip=8pt,after skip=8pt}

\pagestyle{fancy}
\fancyhf{}
\fancyhead[R]{机器学习-第1章}
\fancyhead[L]{机器学习概述 · 练习题}
\fancyfoot[C]{\thepage}
\addtolength{\headheight}{2pt}

\parskip=2pt plus 1pt minus 0.5pt
\linespread{1.35}
\raggedbottom
\widowpenalty=10000
\clubpenalty=10000

\title{\textbf{机器学习\\第1章\quad 机器学习概述}\\[0.5em]
{\large\color{gray}——从图灵的问题到大模型时代，机器学习从哪里来、到哪里去？}}
\date{\today}

\begin{document}
\maketitle

\begin{beginnerbox}[本章题目导览：你将挑战哪些核心问题？]
    机器学习的大厦需要一块坚实的地基：弄清楚它\textbf{是什么、从哪来、能做什么、用什么语言描述}。本章 5 道题覆盖概述全貌：
    \begin{itemize}
        \item \textbf{Q1-01（概念辨析）}：AI、机器学习、深度学习三者的关系——嵌套还是并列？韦恩图怎么画？
        \item \textbf{Q1-02（发展历史）}：六个时代横向对比——每个阶段的核心驱动力、代表成果和局限性是什么？
        \item \textbf{Q1-03（时间线排序）}：把关键里程碑事件按时间排序，建立历史感。
        \item \textbf{Q1-04（术语辨析）}：数据集、特征、标签、模型、训练、推断——6 个基本术语，你真的说得清吗？
        \item \textbf{Q1-05（应用分类）}：给一组实际应用场景，判断它们属于哪种机器学习任务类型。
    \end{itemize}
\end{beginnerbox}

\section{第 1 章\quad 机器学习概述}

在正式学习任何算法之前，先问一个最根本的问题：\textbf{什么是机器学习？}

传统编程的逻辑是：程序员写规则，计算机执行规则，得到输出。但现实中很多问题的规则根本写不出来——你能把"如何识别一只猫"写成一套 \texttt{if-else} 吗？机器学习的回答是：\textbf{不写规则，给数据，让机器自己找规律}。

\begin{metaphorbox}[先建立一个总感觉：机器学习在做什么？]
    \textbf{传统编程}：厨师（程序员）提供菜谱（规则），厨房（计算机）按菜谱做菜（执行），端出菜肴（输出）。\\
    \textbf{机器学习}：给厨房（机器）看一万道菜（数据），让它自己\textbf{总结出做菜的规律}（学习），以后给它新食材（新数据）它能自己做出菜来（预测）。\\
    关键区别：规则是机器\textbf{学出来的}，不是人写的。
\end{metaphorbox}

%% ============================================================
\Needspace{5\baselineskip}
\begin{problembox}[Q1-01\quad AI、机器学习、深度学习：三者关系辨析]
    \textbf{题目描述：}\\
    "人工智能""机器学习""深度学习"这三个词经常被混用，但它们是三个层次不同的概念。

    \vspace{0.3cm}
    \textbf{问题：}
    \begin{enumerate}[(1)]
        \item 用一句话分别定义：人工智能（AI）、机器学习（ML）、深度学习（DL）。
        \item 三者的关系是\textbf{包含}关系还是\textbf{并列}关系？用文字描述其嵌套结构（最大圈到最小圈分别是哪个概念）。
        \item 填写下表，对比三者的核心特征：
        \begin{center}
        \begin{tabular}{p{2.8cm} p{3.8cm} p{3.8cm} p{2.8cm}}
            \toprule
            概念 & 核心方法 & 典型代表技术/系统 & 兴起时间 \\
            \midrule
            人工智能 & \underline{\hspace{3cm}} & 专家系统、博弈程序 & \underline{\hspace{1.5cm}} \\[1.2em]
            机器学习 & 从数据中自动学习规律 & \underline{\hspace{3cm}} & \underline{\hspace{1.5cm}} \\[1.2em]
            深度学习 & \underline{\hspace{3cm}} & \underline{\hspace{3cm}} & 2010年代 \\
            \bottomrule
        \end{tabular}
        \end{center}
        \item 举一个例子，说明某个任务\textbf{属于 AI 但不属于机器学习}（即用规则/逻辑而非数据学习的方式解决）。
    \end{enumerate}
    {\color{gray}\small\textit{来源溯源：[文档第 1 章 1.1--1.2]}}
\end{problembox}

\begin{solutionbox}[参考答案与详细解析]
    \textbf{(1) 一句话定义：}
    \begin{itemize}
        \item \textbf{人工智能（AI）}：让机器具备类似人类智能的技术与研究领域的总称，包括推理、规划、感知、自然语言理解等能力。
        \item \textbf{机器学习（ML）}：AI 的一个子领域，研究如何让机器通过数据自动学习规律，而无需显式编程规则。
        \item \textbf{深度学习（DL）}：机器学习的一个子领域，专指使用多层神经网络（深层架构）从数据中自动提取层次化特征的方法。
    \end{itemize}

    \textbf{(2) 嵌套包含关系：}

    AI（最大圈）$\supset$ ML（中圈）$\supset$ DL（最小圈）。深度学习是机器学习的一种方法，机器学习是人工智能的一种实现途径，人工智能还包括规则系统、专家系统等非机器学习方法。

    \textbf{(3) 对比表格：}
    \begin{center}
    \begin{tabular}{p{2.8cm} p{3.8cm} p{3.8cm} p{2.8cm}}
        \toprule
        概念 & 核心方法 & 典型代表 & 兴起时间 \\
        \midrule
        人工智能 & 规则推理、逻辑符号操作、搜索 & 专家系统、深蓝（国际象棋程序） & 1950年代 \\[0.8em]
        机器学习 & 从数据中自动学习规律 & SVM、决策树、随机森林 & 1980--90年代 \\[0.8em]
        深度学习 & 多层神经网络自动特征提取 & CNN、Transformer、GPT & 2010年代 \\
        \bottomrule
    \end{tabular}
    \end{center}

    \textbf{(4) 属于 AI 但不属于 ML 的例子：}

    \textbf{国际象棋的 Minimax 搜索算法}：通过穷举博弈树、用启发式估值函数评估局面，完全基于规则和搜索，没有从数据中学习任何参数。又如早期的\textbf{专家系统}（如 MYCIN 医疗诊断系统），将领域专家的知识编码为 IF-THEN 规则，靠规则推理而非数据训练。
\end{solutionbox}

%% ============================================================
\Needspace{5\baselineskip}
\begin{problembox}[Q1-02\quad 发展历史：六个时代的核心特征对比]
    \textbf{题目描述：}\\
    机器学习经历了六个主要发展阶段，每个阶段有其特定的核心驱动力、代表成果和历史局限。

    \vspace{0.3cm}
    \textbf{问题：}
    \begin{enumerate}[(1)]
        \item 填写下表，概括各阶段的关键信息：
        \begin{center}
        \begin{tabular}{p{3.5cm} p{2.5cm} p{3.5cm} p{3cm}}
            \toprule
            阶段名称 & 时间范围 & 核心驱动力/代表成果 & 主要局限 \\
            \midrule
            早期探索 & \underline{\hspace{1.5cm}} & 图灵测试、感知机 & \underline{\hspace{2cm}} \\[1.2em]
            知识驱动与专家系统 & 70--80年代 & \underline{\hspace{3cm}} & \underline{\hspace{2cm}} \\[1.2em]
            数据驱动与统计学习 & \underline{\hspace{1.5cm}} & SVM、概率图模型 & \underline{\hspace{2cm}} \\[1.2em]
            深度学习崛起 & 21世纪初 & \underline{\hspace{3cm}} & \underline{\hspace{2cm}} \\[1.2em]
            大模型时代 & \underline{\hspace{1.5cm}} & \underline{\hspace{3cm}} & 对齐、安全、能耗 \\
            \bottomrule
        \end{tabular}
        \end{center}
        \item 为什么 20 世纪 70--80 年代专家系统盛行，却在 80 年代末迎来"AI 寒冬"？知识瓶颈的核心矛盾是什么？
        \item 深度学习在 2010 年代崛起的三个关键条件是什么？（提示：不只是算法层面的突破）
        \item 大模型（如 GPT 系列）与传统机器学习模型在\textbf{规模}和\textbf{训练范式}上有哪两个本质区别？
    \end{enumerate}
    {\color{gray}\small\textit{来源溯源：[文档第 1 章 1.3.1--1.3.5]}}
\end{problembox}

\begin{solutionbox}[参考答案与详细解析]
    \textbf{(1) 六阶段对比表：}
    \begin{center}
    \begin{tabular}{p{3.5cm} p{2.5cm} p{3.5cm} p{3cm}}
        \toprule
        阶段 & 时间 & 核心成果 & 主要局限 \\
        \midrule
        早期探索 & 50--70年代 & 图灵测试、感知机、LISP & 计算力不足，只能处理玩具问题 \\[0.8em]
        知识驱动/专家系统 & 70--80年代 & MYCIN、DENDRAL 等专家系统 & 知识获取瓶颈，无法处理不确定性 \\[0.8em]
        数据驱动/统计学习 & 80年代--21世纪初 & SVM、贝叶斯网络、AdaBoost & 特征工程依赖人工，浅层模型表达受限 \\[0.8em]
        深度学习崛起 & 2006--2015 & AlexNet、Word2Vec、GAN & 需要大量标注数据，可解释性差 \\[0.8em]
        大模型时代 & 2017--至今 & GPT、BERT、Stable Diffusion & 对齐问题、安全风险、巨大能耗 \\
        \bottomrule
    \end{tabular}
    \end{center}

    \textbf{(2) 专家系统的"知识瓶颈"：}

    专家系统需要将人类专家的知识逐条手工编码为规则，面临三大矛盾：\textbf{知识难以穷举}（现实世界的规则无限多）；\textbf{知识难以更新}（规则冲突时难以自动调整）；\textbf{无法处理不确定性}（规则系统缺乏对模糊/噪声情况的处理能力）。这使得专家系统在封闭领域有效但无法泛化，维护成本极高，最终导致资金和信心的双重崩溃。

    \textbf{(3) 深度学习崛起的三个条件：}
    \begin{itemize}
        \item \textbf{大数据}：互联网的爆炸式增长提供了海量带标注的训练数据（如 ImageNet 的 120 万张图片）；
        \item \textbf{算力提升}：GPU 的并行计算能力使深层网络训练在合理时间内成为可能；
        \item \textbf{算法突破}：ReLU 激活函数、Dropout 正则化、Batch Normalization 等技巧解决了梯度消失问题。
    \end{itemize}

    \textbf{(4) 大模型 vs 传统 ML 的两个本质区别：}
    \begin{itemize}
        \item \textbf{规模}：参数量从百万级（传统 ML）跃升至千亿级（GPT-3 1750 亿参数），涌现出传统模型不具备的"涌现能力"（in-context learning、chain-of-thought 等）；
        \item \textbf{训练范式}：传统 ML 针对特定任务训练专用模型；大模型采用"预训练 + 微调/提示"范式——先在超大规模无标注数据上预训练通用表示，再以极少标注数据适配下游任务。
    \end{itemize}
\end{solutionbox}

%% ============================================================
\Needspace{5\baselineskip}
\begin{problembox}[Q1-03\quad 历史里程碑时间线排序]
    \textbf{题目描述：}\\
    下面列出了 10 个机器学习领域的重要里程碑事件，请按时间先后排序（最早到最晚），并在每个事件旁标注大致年份。

    \vspace{0.2cm}
    \textbf{事件列表（乱序）：}
    \begin{enumerate}[A.]
        \item Transformer 架构论文《Attention is All You Need》发表
        \item 图灵发表论文《计算机器与智能》，提出图灵测试
        \item AlexNet 在 ImageNet 竞赛中以大幅优势夺冠，引爆深度学习热潮
        \item Minsky 和 Papert 出版《感知机》，指出单层感知机局限性，引发第一次 AI 寒冬
        \item Rosenblatt 发明感知机（Perceptron）
        \item Hinton 等人提出深度信念网络，深度学习概念复兴
        \item ChatGPT 发布，大语言模型进入大众视野
        \item MYCIN 专家系统用于医疗诊断，专家系统达到鼎盛
        \item Vapnik 提出支持向量机（SVM）
        \item 反向传播算法被重新发现并推广，神经网络研究复苏
    \end{enumerate}

    \vspace{0.2cm}
    \textbf{问题：}
    \begin{enumerate}[(1)]
        \item 将上述事件按时间顺序排列（写出字母序列），并在每个字母后括号内填写大致年份。
        \item 指出哪两个事件之间发生了"AI 寒冬"，并说明原因。
        \item 哪个事件标志着深度学习时代的\textbf{正式崛起}？为什么以此为标志？
    \end{enumerate}
    {\color{gray}\small\textit{来源溯源：[文档第 1 章 1.3.1--1.3.5]}}
\end{problembox}

\begin{solutionbox}[参考答案与详细解析]
    \textbf{(1) 时间顺序排列：}

    \textbf{B（1950）→ E（1957）→ D（1969）→ H（1974）→ J（1986）→ I（1995）→ F（2006）→ C（2012）→ A（2017）→ G（2022）}

    \textbf{(2) AI 寒冬发生在 D 和 J 之间（1969--1986）：}

    Minsky 等人的《感知机》从数学上证明了单层感知机的局限（无法解决 XOR 等非线性问题），导致神经网络研究经费大幅削减。直到 1986 年反向传播算法重新流行，多层网络的训练问题才得到解决。

    \textbf{(3) 事件 C（2012 年 AlexNet）标志深度学习崛起：}

    AlexNet 在 ImageNet 竞赛中将 Top-5 错误率从 26\% 降低到 15.3\%，远超第二名，差距之大震惊学界。它首次将深层卷积网络、GPU 加速、ReLU 激活、Dropout 正则化结合，\textbf{用实验成果而非理论论证}，彻底说服了整个计算机视觉领域，引发深度学习的全面爆发。
\end{solutionbox}

%% ============================================================
\Needspace{5\baselineskip}
\begin{problembox}[Q1-04\quad 基本术语辨析：数据集、特征、标签……]
    \textbf{题目描述：}\\
    机器学习有一套专用术语，同一个概念在不同语境下有不同名称，初学者常常混淆。

    \vspace{0.3cm}
    \textbf{问题：}
    \begin{enumerate}[(1)]
        \item 填写下表，给出每个术语的定义，并用"预测房价"这一场景举例说明：

        \begin{center}
        \begin{tabular}{p{2.5cm} p{4.5cm} p{4.5cm}}
            \toprule
            术语 & 定义 & "预测房价"场景中的对应含义 \\
            \midrule
            样本（Sample） & \underline{\hspace{4cm}} & 某一套具体的房子的信息 \\[1.2em]
            特征（Feature） & \underline{\hspace{4cm}} & \underline{\hspace{4cm}} \\[1.2em]
            标签（Label） & \underline{\hspace{4cm}} & \underline{\hspace{4cm}} \\[1.2em]
            训练集 & \underline{\hspace{4cm}} & \underline{\hspace{4cm}} \\[1.2em]
            测试集 & \underline{\hspace{4cm}} & \underline{\hspace{4cm}} \\[1.2em]
            模型（Model） & \underline{\hspace{4cm}} & 根据房屋特征预测价格的函数 \\
            \bottomrule
        \end{tabular}
        \end{center}

        \item "训练"（Training）和"推断/预测"（Inference）有什么区别？各自在什么阶段发生？
        \item 什么叫"泛化能力"（Generalization）？为什么我们不用训练集上的表现来评估模型好坏？
        \item 在机器学习中，"特征"有时也叫"属性"或"输入变量"，"标签"有时也叫"目标变量"或"输出变量"。结合这些别名，说明特征和标签在数据表格中各自对应什么列？
    \end{enumerate}
    {\color{gray}\small\textit{来源溯源：[文档第 1 章 1.4 基本术语]}}
\end{problembox}

\begin{solutionbox}[参考答案与详细解析]
    \textbf{(1) 术语表格：}
    \begin{center}
    \begin{tabular}{p{2.5cm} p{4.5cm} p{4.5cm}}
        \toprule
        术语 & 定义 & "预测房价"示例 \\
        \midrule
        样本 & 数据集中的一条记录，代表一个观测对象 & 某一套具体的房子的所有信息（一行数据） \\[0.8em]
        特征 & 描述样本属性的输入变量（可以有多个） & 面积、楼层、朝向、地铁距离、建成年份等 \\[0.8em]
        标签 & 样本对应的真实输出/答案，是模型要预测的目标 & 该房子的实际成交价格 \\[0.8em]
        训练集 & 用于训练模型（调整参数）的数据集 & 已知价格的历史房屋数据，用来拟合模型 \\[0.8em]
        测试集 & 训练结束后用于评估模型泛化能力的数据集，训练时不可见 & 另一批房屋数据，用来检验预测准不准 \\[0.8em]
        模型 & 从特征到标签的映射函数，通过训练从数据中学习得到 & 根据房屋特征预测价格的函数 \\
        \bottomrule
    \end{tabular}
    \end{center}

    \textbf{(2) 训练 vs 推断：}
    \begin{itemize}
        \item \textbf{训练}：使用训练集数据，通过优化算法（如梯度下降）反复调整模型参数，使模型在训练数据上的预测误差逐渐减小。发生在\textbf{建模阶段}。
        \item \textbf{推断/预测}：训练完成后，将新的未见过的输入特征送入固定参数的模型，得到预测结果。发生在\textbf{部署/应用阶段}，此时模型参数不再更新。
    \end{itemize}

    \textbf{(3) 泛化能力：}

    泛化能力指模型对\textbf{未见过的新数据}做出正确预测的能力。训练集上的表现无法衡量泛化——因为模型可以通过"死记硬背"训练数据达到零误差（过拟合），而对新数据的预测却很差。只有在训练时\textbf{完全没有接触过}的测试集上测出的误差，才能真实反映模型的实用价值。

    \textbf{(4) 特征与标签在表格中的位置：}

    在数据表格中，特征（属性/输入变量）对应除目标列外的\textbf{所有其他列}（通常是前几列）；标签（目标变量/输出变量）对应\textbf{最后一列}（或单独指定的目标列）。每一行是一个样本，每一列是一个变量。
\end{solutionbox}

%% ============================================================
\Needspace{5\baselineskip}
\begin{problembox}[Q1-05\quad 应用场景分类：判断机器学习任务类型]
    \textbf{题目描述：}\\
    机器学习根据数据是否有标签、输出类型等维度，分为监督学习、无监督学习、强化学习等类型。下面给出若干实际应用场景，请判断每个场景最适合用哪类机器学习方法解决，并说明理由。

    \vspace{0.2cm}
    \textbf{应用场景列表：}
    \begin{enumerate}[A.]
        \item 银行根据客户的历史交易记录，预测某笔交易是否为欺诈
        \item 电商平台根据用户浏览行为（无购买标签），自动将用户分成若干消费群体
        \item 游戏 AI 通过反复自我对弈，学习如何在围棋比赛中取得最高得分
        \item 医疗系统根据患者的症状描述文本，自动将其分类到对应科室（有历史挂号数据作为标注）
        \item 将一张 $1024\times1024$ 的高清照片压缩到更小尺寸，同时尽量保留主要内容（无标签）
        \item 自动驾驶汽车通过与环境的持续互动，学习在不同路况下的驾驶策略
    \end{enumerate}

    \vspace{0.2cm}
    \textbf{问题：}
    \begin{enumerate}[(1)]
        \item 对每个场景，填写：任务类型（监督/无监督/强化学习）、具体子类型（分类/回归/聚类/降维/序列决策等）、判断理由。
        \item 监督学习中，"分类"和"回归"的根本区别是什么？各举两个现实应用例子。
        \item 机器学习还有一种重要范式——\textbf{自监督学习}（Self-supervised Learning），它与无监督学习有什么区别？GPT 的预训练属于哪种范式？
    \end{enumerate}
    {\color{gray}\small\textit{来源溯源：[文档第 1 章 1.3.6 机器学习应用领域]}}
\end{problembox}

\begin{solutionbox}[参考答案与详细解析]
    \textbf{(1) 场景分类：}
    \begin{itemize}
        \item \textbf{A}：\textbf{监督学习——二分类}。有历史交易记录（特征）和是否欺诈的标注（标签），目标是预测"是/否"这一离散类别。
        \item \textbf{B}：\textbf{无监督学习——聚类}。没有预定义的群体标签，仅凭用户行为数据自动发现内在分组结构。
        \item \textbf{C}：\textbf{强化学习——序列决策}。智能体通过与环境交互，根据胜负信号（奖励）学习最优策略，没有"正确答案"标签。
        \item \textbf{D}：\textbf{监督学习——多分类}。有历史挂号数据作为标注（特征=症状文本，标签=科室），目标是从多个类别中选一个。
        \item \textbf{E}：\textbf{无监督学习——降维}。无标签，目标是在保留主要信息的前提下减少数据维度（如 PCA 降维）。
        \item \textbf{F}：\textbf{强化学习——连续控制}。自动驾驶需要在动态环境中持续做决策，通过与环境交互和奖励信号学习策略。
    \end{itemize}

    \textbf{(2) 分类 vs 回归：}

    根本区别在于\textbf{输出类型}：分类的输出是\textbf{离散的类别标签}（有限个选项），回归的输出是\textbf{连续的数值}。
    \begin{itemize}
        \item \textbf{分类示例}：垃圾邮件检测（是/否）；图像识别（猫/狗/鸟）。
        \item \textbf{回归示例}：房价预测（具体金额）；气温预测（具体度数）。
    \end{itemize}

    \textbf{(3) 自监督学习 vs 无监督学习：}

    \begin{keybox}[自监督学习的核心思想]
        无监督学习不使用任何标签，纯粹发现数据结构；\textbf{自监督学习}则从数据本身\textbf{自动构造监督信号}——比如遮住句子中的某些词，让模型预测被遮住的词（BERT 的 MLM 任务），或用前文预测下一个词（GPT 的 CLM 任务）。标签是从数据中自动生成的，不需要人工标注。GPT 的预训练属于\textbf{自监督学习}（语言建模任务）。
    \end{keybox}
\end{solutionbox}

\begin{appbox}[现实应用场景]
    \textbf{机器学习的六大应用领域一览：}\\
    计算机视觉（图像分类、目标检测、人脸识别）、自然语言处理（机器翻译、情感分析、问答系统）、语音识别与合成、推荐系统（电商、短视频、音乐）、医疗健康（辅助诊断、药物研发、基因组学）、金融科技（风控、量化交易、反欺诈）。这六大领域贡献了机器学习绝大多数的商业价值，也是就业最集中的方向。
\end{appbox}

\section{本章小结}
\begin{itemize}
    \item \textbf{ML 是 AI 的子集，DL 是 ML 的子集}：三层嵌套，不要混用。
    \item \textbf{六个时代的主旋律}：规则→知识→统计→深度→大模型，每次转变都是"数据量和算力突破了瓶颈"。
    \item \textbf{基本术语是语言基础}：样本/特征/标签/训练集/测试集/模型，是所有后续章节的共同语言。
    \item \textbf{任务类型决定方法选择}：有标签→监督学习；无标签→无监督；序列决策→强化学习。
\end{itemize}
\end{document}
